% !PRINTED 585
\weberpage{585}
%<<CONT
daher nur das Product aller von einander verschiedenen unter den conjugirten Idealen nimmt, so braucht dies keineswegs ein Ideal in $R$ zu sein.

Was hier von den Normen ausgef\"uhrt ist, gilt auch von den anderen symmetrischen Functionen, insbesondere von den Spuren. Wir nennen die Summe
\weq{6}{\mathfrak{S}_R(\omega) = \omega_1 + \omega_2 + \cdots + \omega_n}
die \gesperrt{Partialspur (oder Relativspur) von $\omega$ in Bezug auf} $R$. Sie ist eine Zahl oder ein Functional in $R$, und wir erhalten f\"ur die Totalspur
\weq{7}{S_\Omega(\omega) = S_R\,\mathfrak{S}_R(\omega),}
worin die Bedeutung der Zeichen unmittelbar verst\"andlich ist.

Eine besonders wichtige Beziehung ergiebt sich aber durch die Betrachtung der Grundideale.

Es sei $\tau$ eine Basisform von $\Omega$ (also eine Linearform von $mn$ Variablen, \S.\,146) und
\[
F(t) = N_\Omega(t - \tau)
\]
die irreducible Function $mn^{\text{ten}}$ Grades (im absoluten Rationalit\"atsbereiche), deren Wurzel $\tau$ ist. Dann haben wir im \S.\,158 das durch das Functional $F'(\tau)$ definirte Ideal als das \gesperrt{Grundideal} des K\"orpers $\Omega$ bezeichnet.

Ist nun $\sigma$ eine Basisform von $R$ und
\[
\psi(t) = N_R(t - \sigma),
\]
so ist durch $\psi'(\sigma)$ das Grundideal des K\"orpers $R$ definirt.

Aus der Function $F(t)$ spaltet sich aber im K\"orper $R$ ein irreducibler Factor $n^{\text{ten}}$ Grades ab,
\[
f(t) = \mathfrak{N}_R(t - \tau),
\]
der f\"ur $t = \tau$ verschwindet. Wir definiren durch $f'(\tau)$ das \gesperrt{Partial-Grundideal von $\Omega$ in Bezug auf} $R$, und beweisen nun den Satz:
\weq{8}{F'(\tau) = \varepsilon\, f'(\tau)\, \psi'(\sigma),}
worin $\varepsilon$ eine Einheit ist.

Um ihn zu beweisen, setzen wir nach der im \S.\,158 angewandten Bezeichnung
\weq{9}{\frac{F(t)}{t - \tau} = t^{mn-1}F_0(\tau) + t^{mn-2}F_1(\tau) + \cdots + F_{mn-1}(\tau),}
\weq{10}{\frac{f(t)}{t - \tau} = t^{n-1}f_0(\tau) + t^{n-2}f_1(\tau) + \cdots + f_{n-1}(\tau),}
\weq{11}{\frac{\psi(t)}{t - \sigma} = t^{m-1}\psi_0(\sigma) + t^{m-2}\psi_1(\sigma) + \cdots + \psi_{m-1}(\sigma).}
